\section{Scope, Method, and Qualifications}

\subsection{Reader, question, and claim boundary}

This is a discursive article in faith and theology for an intellectually serious reader---Christian or considering the Christian claim---who can follow a metaphysical distinction but is not presumed to know scholastic terminology. Its governing question: does the concrete particularity of the biblical economy---one chosen people, one incarnate life, instituted material sacraments, personal address---degrade the universality proper to God, as a long objection tradition running from Celsus through Spinoza, Lessing, and Kant to contemporary religious pluralism contends? Its governing claim: the objection rests on a false premise that generality is the divine mode; in the classical account, only particulars exist concretely, abstraction is the finite knower's limitation, individuality is a positive perfection whose summit in creation is the person, God knows and sustains every singular directly, and the biblical particulars are consistently ordered to universal scope---election for all nations, one assumed body grounding union with every human being, repeatable concrete signs carrying the once-for-all event outward.

The argument establishes coherence and fittingness, not occurrence. It does not prove the Incarnation or Resurrection from metaphysics, adjudicate among claimed revelations, resolve historical questions such as the Quirinius census date or the Lucan synchronisms, or supply the reader's act of faith. Fittingness arguments (\emph{convenientia}) are used only retrospectively: they show why an event believed on other grounds is worthy of God, never that it had to happen.

\subsection{Terminology fixed for this article}

\begin{itemize}
\item \textbf{Particular} means concrete, singular, historically located: this person, people, body, act, sign. It is predicated of creatures, divine actions, and the assumed humanity of Christ---never of God as though he were one member of a class.
\item \textbf{Universal} is used in two senses kept explicit in the argument: the abstract universal (a common nature as object of thought) and the universal cause (God as giver of being to every creature). The article's central distinction is that the second does not entail the first as a divine mode of relating to creatures.
\item \textbf{Individuation} names whatever makes a particular the particular it is; the article expounds the Thomist answer (designated matter) and reports the Scotist answer (haecceity) without adjudicating between them.
\item \textbf{Person} is used in Boethius's sense as expounded by Aquinas (a subsistent individual of a rational nature), and is predicated of God only analogically, as the cited article itself requires.
\item \textbf{Scandal of particularity} is a received modern label for the objection, not a scriptural or magisterial category.
\item \textbf{Flight into abstraction} names the article's diagnosis of the preference for a non-landing God; it is an editorial characterization, not a canonical classification of any thinker.
\item \textbf{Election} means gratuitous divine choice conferring vocation and mission, expressly not antecedent merit or exclusive divine concern.
\item \textbf{Pluralist hypothesis} means the position summarized at the reported ceiling in the contemporary section (one ultimate Reality, many equally valid culturally shaped responses); it is engaged as a philosophical-theological thesis, not as a description of any named person's faith.
\end{itemize}

\subsection{Source hierarchy and authority classes}

Scripture governs the pattern claims (election's gratuity and outward ordering through the patriarchal narrowing, the prophets' universalism-through-particularity, the Incarnation and its genealogies, the risen Christ's particularity, the Church's one baptism); brief clauses are quoted from the public-domain Douay--Rheims (Challoner) translation as an identified working witness, without any claim that it is a critical text. Conciliar, catechetical, and dicasterial texts (\emph{Gaudium et spes} 22; \emph{Catechism of the Catholic Church} 1084--1085 and 1145--1152; \emph{Dominus Iesus} 2, 4, 6, 9--10, 13--15, 20--22) govern the universality of Christ's union with humanity, the sacramental economy, and the answer to pluralism, each at its own authority level; \emph{Dominus Iesus} is cited as a dicasterial declaration reiterating doctrine taught in the sources it cites. Patristic witnesses (Origen answering Celsus; Irenaeus; Athanasius; Leo) show that the particularity objection was met, not evaded, in the tradition; breadth of citation is illustrative, not an authority count. Aquinas supplies the metaphysical framework (knowledge of singulars, providence over individuals, God outside every genus and innermostly present to each thing, individuation by designated matter, the person as summit of created perfection, the individual term of the assumption, sensible and instituted signs, the sacraments' conformity to the Incarnation); these are school positions adopted and argued, not dogmas. Celsus (as quoted by Origen), Spinoza, Lessing, Kant, and the pluralist hypothesis are objection sources, quoted or reported to state the strongest opposing case, with their governing premises identified rather than conceded. Secondary encyclopedia witnesses (two named Stanford Encyclopedia of Philosophy entries) are used only at stated reported ceilings. Everything else---the ordering of the argument, the inversion of the abstraction premise, the reading of the genealogies and of Deuteronomy 10's juxtaposition, the letter and doctor and form-letter analogies, the identification of the pluralist hypothesis's Kantian architecture, the reading of Rev 2:17 as particularity perfected---is project synthesis.

\subsection{Material qualifications}

\begin{enumerate}
\item God is not a particular alongside others; predicating choice, address, and personality of God does not place him in a genus (ST I, q.~3, a.~5 and I, q.~29, a.~3 are cited to that effect, the latter with its own analogical qualification preserved).
\item The moderate-realist account of universals, the Thomistic analysis of abstraction, and the doctrine of individuation by designated matter are adopted school positions; the Scotist alternative is reported at a named-witness ceiling and the dispute is not adjudicated; the theological claims of the later sections do not stand or fall with any one technical formulation.
\item The angelology of ST I, q.~50, a.~4 is cited only to show the individuation principle applied consistently; nothing depends on it.
\item The essence--existence distinction, participated being, and analogical predication are argued in companion studies of this collection and are borrowed here, not re-argued; the section using them states a weaker fallback for readers who reject the framework.
\item Divine knowledge of singulars and particular providence follow the classical Thomistic account; no inference is licensed from particular suffering to particular guilt.
\item Election's outward ordering does not by itself identify which claimed revelation is true; the article says so explicitly and argues coherence only. Israel's election remains, in Christian teaching, irrevocable; nothing here treats Israel as a discarded instrument, and the Israel--Church relation is not adjudicated. Deuteronomy 7 is cited for vv.~6--8 only; the surrounding herem commands raise questions outside this article's scope and receive no endorsement, express or implied.
\item Genesis 25:23 is cited for the pattern of choice antecedent to merit; the dogmatic questions opened by Romans 9's use of it are outside scope. The readings of Matthew's four women, of Luke's genealogy, and of Deuteronomy 10:14--15's juxtaposition are labeled interpretations.
\item Athanasius's cosmic-Word argument is directed \emph{ad hominem} at interlocutors who granted the Word's presence in the whole; the hypostatic union differs in kind from the Word's union with creation, as the cited edition's own annotation notes. Origen's replies to Celsus are used for the two points quoted (providence beyond the whole; non-topical descent; the sufficiency of the one Word), not adopted wholesale.
\item ``United Himself in some fashion with every man'' retains the council's qualifier; universal salvific scope is not converted into a claim that all are in fact saved.
\item Sacramental making-present is not repetition of the historical event; the paragraphs on assurance do not reduce grace to mechanism or dispense with faith and disposition; sacramental validity, matter and form, and the exposition or defense of Eucharistic doctrine (including transubstantiation) are outside scope---ST III, qq.~75--76 are quoted for the fittingness reasons and the whole-in-each-part structure only, and the real presence is stated as the Church's teaching, not argued philosophically here.
\item The Lessing quotations are given in German from the Lachmann--Muncker \emph{S\"amtliche Schriften}, 3rd ed., vol.~13 (Stuttgart, 1897), pp.~5 and 7, in that edition's eighteenth-century orthography; the English renderings are the article's own labeled working translations, and no public-domain English translation was identified. The 1777 Brunswick first printing, which that edition's editor names as the only textually authoritative witness, was not itself inspected, and this remaining limit is recorded here and in the research record.
\item Kant is quoted from Semple's 1838 public-domain translation as an identified working witness; the German was not collated, and no claim finer than the quoted argumentative structure is made. Kant's comparative characterizations of other religions are expressly not reproduced or endorsed.
\item Hick and the surrounding debate are engaged entirely at a stated reported ceiling on a named SEP witness with its revision date; no Hick text was collated, and no claim about his position finer than that witness supports is made. The same holds for the Scotus material and its SEP witness.
\item \emph{Dominus Iesus} is quoted exactly from the official English text at the Vatican website; its teaching on the situation of non-Christians is reported at its own loci without extension, and this article makes no claim about the salvation or standing of any individual or named community.
\item Celsus survives only in Origen's quotation; attribution is to Celsus as reported, and the ancient polemic's contempt for the Jewish people, where quoted to state the objection, is reported and rejected, not endorsed.
\item The closing section's claims about love and the white stone are theological reading and synthesis, not exegetical assertions about authorial intention.
\end{enumerate}

\subsection{Coverage, currentness, translations, rights, and review}

No liturgical rite, civil or canonical jurisdiction, or mutable discipline governs any conclusion; no legal claims are made. Online witnesses were checked on 24 July 2026 (first edition) and 25 July 2026 (deepened edition): Douay--Rheims chapter pages; New Advent transcriptions of the Ante-Nicene and Nicene and Post-Nicene Fathers translations (including \emph{Contra Celsum} books IV--VI) and of the English Dominican \emph{Summa} (including I, qq.~8, 29, 50; III, qq.~60--61, 75--76); a full-text NPNF copy of Robertson's Athanasius; the Corpus Thomisticum Latin for ST III, q.~1, a.~1 and for \emph{De ente et essentia}; the Vatican website's English \emph{Catechism}, \emph{Gaudium et spes}, and \emph{Dominus Iesus}; the Project Gutenberg Elwes translation of Spinoza; the Internet Archive scans of Semple's 1838 Kant and of Lachmann--Muncker's Lessing, vol.~13; and the named Stanford Encyclopedia witnesses with their revision dates recorded. Quotations are brief and argument-driven. Scripture, official ecclesiastical texts, and third-party translations retain their own status and are not offered under the project's CC BY 4.0 license; the Stanford Encyclopedia entries are used as checked secondary narratives and quoted only in short clauses with attribution. This is an internally source-audited working study by an AI-assisted process; the 25 July 2026 deepening (objection genealogy at depth with the Kant subsection, individuation and person material, expanded scriptural and sacramental sections, the contemporary pluralism section with \emph{Dominus Iesus}, and the synopsis table) passed through the same audit; it has received no independent philosophical, biblical, patristic, dogmatic, literary, or ecclesiastical review and claims no ecclesiastical approval.
